\section{Spline Sufficiency}

Finite measurement gives the grammar anchors before it gives the grammar a
curve. An anchor is a committed reading: an event, a mark, a boundary value,
a recorded position, a calibration point. Between anchors, the record is
silent. The silence is not empty permission. The grammar may complete the
record only in ways that remain answerable to the anchors and to the
residual tests already established. A completion that inserts a hidden
distinction has made a new claim, and the new claim must pay for itself.

A spline is the disciplined form of this completion. It does not say that a
continuous object was present behind the finite record all along. It says
that, once the anchors are fixed, the grammar may assign intermediate
values only by the least strained rule compatible with those anchors. The
intermediate values are consequences of the completion. They are not new
measurements. If a later refinement records a new anchor, the completion
must answer to it. Until then, the spline is the permitted smooth shadow of
the finite ledger.

The cubic case is the first case rich enough to carry the needed balance.
A line can connect two anchors but has no bending to account for. A
quadratic can bend but cannot generally balance neighboring pieces without
leaving a derivative debt at the knot. A cubic patch has enough freedom to
match value, slope, and bending moment across an anchor. It can carry the
position, the first slip, and the second slip without inventing a new
interior fact.

This is why the chapter speaks of a closed balanced cubic spline. Balanced
means that adjacent patches agree at their shared anchor in the readings the
gauge is allowed to test. The value agrees, so the path does not jump. The
first reading agrees, so the path does not hide a corner. The bending
reading agrees, so the path does not hide a curvature debt. Closed means
that the boundary of the chain has no unspent residue: the completion
returns to the committed boundary ledger without adding a final slip.

Once these requirements hold, the first residual has nowhere to live. Any
supported variation inside a patch is already governed by the cubic
balance. Any variation at a knot is answered by the matching conditions
from the two neighboring patches. Any variation at the boundary is
forbidden by the committed endpoint or cancelled by closedness. The grammar
therefore reads residue zero. Spline sufficiency is exactly this claim: a
closed balanced cubic completion is sufficient to force stationarity.

The claim is modest in the right way. It does not say that every stationary
path must be presented as a cubic spline. It says that this finite spline
grammar supplies enough structure to make residual vanishing unavoidable.
The anchors carry the finite facts; the cubic patches carry the least
completion; the knot balances carry the local Euler--Lagrange reading; the
closure carries the boundary debt to zero. Nothing else is needed for the
first residual to vanish.

The distinction between precision and accuracy is useful here. A spline can
make a value at an unmeasured point precise in the sense that the completion
identifies a unique value there. That precision is a grammatical
consequence of the anchors and the completion rule. It is not accuracy in
the empirical sense. If a later observation places a new anchor elsewhere,
the ledger changes and the completion must be recomputed. The grammar
therefore permits precise intermediate readings while forbidding them from
masquerading as independently recorded facts.

This is also why hidden curvature cannot be admitted for free. Suppose two
readers refine the same interval between the same anchors. If a hidden bend
were large enough to be distinguishable, refinement would expose a new
event. If no refinement exposes such an event, the permitted completion
must collapse toward the same balanced cubic patch. Agreement of value,
slope, and bending moment at shared anchors is not an aesthetic condition.
It is the grammar's way of forbidding a contradiction at the place where
two finite records meet.

Spline sufficiency therefore turns the stationary reader into a constructive
finite shape. The grammar has not assumed a continuum; it has permitted a
completion. It has not assumed smoothness; it has forced enough smoothness
to prevent unrecorded residue. It has not assumed a law of motion; it has
identified the closure condition under which all allowed first variations
read zero. The path is stationary because the balanced spline leaves no
first-order imbalance for the gauge to detect.

The next problem is scale. A single spline over a finite set of anchors
gives a local completion, but measurement also refines. Partitions split,
supports shrink, and residuals are tested on finer ledgers. The grammar
must explain how these finite completions are squeezed toward a completed
reading without treating the completed reading as primitive.
